\section*{Глава 3. Разработка веб-приложения}
\addcontentsline{toc}{section}{Глава 3. Разработка веб-приложения}
\stepcounter{section}

\subsection{Разработка базовой структуры приложения и шаблонов страниц}

В соответствии с выбранным стеком архитектура разрабатываемого программного комплекса построена по принципу разделения клиентской части (SPA на React) и серверного интерфейса (REST API на FastAPI).

Исходный код системы организован в рамках единого репозитория с четким разграничением уровней ответственности:
\begin{itemize}
    \item \texttt{backend/} --- серверная часть на языке Python:
    \begin{itemize}
        \item \texttt{app/core/} --- глобальные параметры конфигурации, настройки безопасности и подключение к СУБД;
        \item \texttt{app/models/} --- декларативные ORM-модели базы данных (SQLAlchemy 2.0);
        \item \texttt{app/schemas/} --- схемы валидации и сериализации входящих и исходящих данных (Pydantic v2);
        \item \texttt{app/services/} --- слой бизнес-логики: адаптеры внешних CI/CD API (GitHub, GitLab), сервис шифрования токенов и модуль диспетчеризации;
        \item \texttt{app/api/v1/} --- модули маршрутизации HTTP-запросов (контроллеры эндпоинтов).
    \end{itemize}
    \item \texttt{frontend/} --- клиентское приложение на React:
    \begin{itemize}
        \item \texttt{src/components/} --- атомарные компоненты интерфейса (кнопки, индикаторы статусов, модальные окна, терминал логов);
        \item \texttt{src/pages/} --- экранные формы приложения (страницы авторизации, главного дашборда, настроек интеграций, детальной аналитики сборок);
        \item \texttt{src/services/} --- клиенты взаимодействия с API на базе Axios с перехватчиками (interceptors) сетевых ответов;
        \item \texttt{src/context/} --- состояние приложения (глобальный контекст авторизации и профиля пользователя).
    \end{itemize}
\end{itemize}

Маршрутизация на стороне клиента реализована с использованием библиотеки \texttt{react-router-dom}. Все защищенные экраны обернуты в компонент \texttt{ProtectedRoute}, который проверяет валидность токена сессии и выполняет перенаправление на страницу входа при его отсутствии.

Базовый каркас интерфейса (\texttt{MainLayout}) включает в себя адаптивную боковую панель навигации (Sidebar), верхнюю строку состояния с информацией о пользователе и центральную контентную область. Стилизация выполнена с использованием утилитарного CSS-фреймворка Tailwind CSS, что обеспечивает консистентный дизайн и адаптивность под различные разрешения экрана.

\subsection{Реализация аутентификации и авторизации пользователей}

Безопасность системы обеспечивается механизмом аутентификации на основе токенов стандарта JWT (JSON Web Token, RFC~7519) в сочетании с алгоритмом хеширования паролей bcrypt.

\subsubsection{Серверная логика безопасности}
При регистрации пользователя введенный пароль не сохраняется в открытом виде, а обрабатывается криптографической функцией хеширования с добавлением соли через библиотеку \texttt{passlib}. Для генерации цифровой подписи токенов применяется алгоритм HMAC-SHA256 с использованием закрытого серверного ключа.

Схема реализации внедрения зависимостей (Dependency Injection) для защиты закрытых эндпоинтов в FastAPI представлена в листинге~\ref{lst:auth_dep}.

\begin{figure}[H]
\centering
\begin{minipage}{0.95\textwidth}
\footnotesize
\begin{verbatim}
from fastapi import Depends, HTTPException, status
from fastapi.security import OAuth2PasswordBearer
from jose import JWTError, jwt
from sqlalchemy.ext.asyncio import AsyncSession

oauth2_scheme = OAuth2PasswordBearer(
    tokenUrl="/api/v1/auth/login"
)

async def get_current_user(
    token: str = Depends(oauth2_scheme),
    db: AsyncSession = Depends(get_db)
) -> User:
    credentials_exception = HTTPException(
        status_code=status.HTTP_401_UNAUTHORIZED,
        detail="Не удалось подтвердить учетные данные",
        headers={"WWW-Authenticate": "Bearer"},
    )
    try:
        payload = jwt.decode(
            token, SECRET_KEY, algorithms=[ALGORITHM]
        )
        user_id: int = payload.get("sub")
        if user_id is None:
            raise credentials_exception
    except JWTError:
        raise credentials_exception

    user = await db.get(User, int(user_id))
    if user is None or not user.is_active:
        raise credentials_exception
    return user
\end{verbatim}
\end{minipage}
\caption{Реализация зависимости проверки JWT-токена в FastAPI}
\label{lst:auth_dep}
\end{figure}

\subsubsection{Клиентская интеграция аутентификации}
На стороне SPA-клиента полученный JWT-токен сохраняется в защищенном хранилище браузера (\texttt{localStorage}). Для взаимодействия с сервером настроен экземпляр Axios, который автоматически подставляет заголовок \texttt{Authorization: Bearer <token>} во все исходящие запросы. При получении HTTP-ответа со статусом 401 (Unauthorized) перехватчик сбрасывает сессию и перенаправляет пользователя на форму входа.

\subsection{Реализация CRUD-интерфейса для управления пайплайнами и этапами выполнения}

Для организации мониторинга пользователь настраивает подключение к целевым Git-платформам и выбирает репозитории для отслеживания.

\subsubsection{Шифрование конфиденциальных данных}
Персональные токены доступа (GitHub PAT, GitLab Token) представляют собой критически важные учетные данные. Сохранение их в открытом виде недопустимо. В модуле \texttt{security.py} реализован алгоритм симметричного шифрования Fernet (на базе AES-128 в режиме CBC с аутентификацией HMAC), представленный в листинге~\ref{lst:fernet_crypto}.

\begin{figure}[H]
\centering
\begin{minipage}{0.95\textwidth}
\footnotesize
\begin{verbatim}
from cryptography.fernet import Fernet
from app.core.config import settings

cipher = Fernet(settings.ENCRYPTION_KEY.encode())

def encrypt_token(plain_token: str) -> str:
    return cipher.encrypt(plain_token.encode()).decode()

def decrypt_token(encrypted_token: str) -> str:
    return cipher.decrypt(encrypted_token.encode()).decode()
\end{verbatim}
\end{minipage}
\caption{Функции шифрования и дешифрования токенов внешних провайдеров}
\label{lst:fernet_crypto}
\end{figure}

\subsubsection{Спецификация маршрутов REST API}
Для управления сущностями разработаны программные эндпоинты, спецификация которых приведена в таблице~\ref{tab:api_endpoints}.

\begin{table}[H]
    \centering
    \caption{Спецификация эндпоинтов управления (префикс \texttt{/api/v1})}
    \label{tab:api_endpoints}
    \small
    \resizebox{\textwidth}{!}{%
    \begin{tabular}{|l|l|p{8.5cm}|}
        \hline
        \textbf{Метод} & \textbf{Маршрут} & \textbf{Назначение операции} \\ \hline
        \texttt{POST} & \texttt{/auth/login} & Аутентификация пользователя и выпуск JWT \\ \hline
        \texttt{GET} & \texttt{/providers} & Получение списка подключенных провайдеров \\ \hline
        \texttt{POST} & \texttt{/providers} & Подключение провайдера (с шифрованием токена) \\ \hline
        \texttt{DELETE} & \texttt{/providers/\{id\}} & Удаление интеграции и связанных репозиториев \\ \hline
        \texttt{GET} & \texttt{/repositories} & Список отслеживаемых репозиториев с фильтрацией \\ \hline
        \texttt{POST} & \texttt{/repositories/sync} & Синхронизация списка репозиториев из внешнего API \\ \hline
        \texttt{PATCH} & \texttt{/repositories/\{id\}} & Переключение флага мониторинга (\texttt{is\_monitored}) \\ \hline
    \end{tabular}%
    }
\end{table}

На клиентской стороне разработан интерфейс настроек (\texttt{SettingsPage}), позволяющий в интерактивном режиме вводить токены провайдеров, проверять соединение с внешним API и включать отслеживание необходимых проектов с помощью переключателей (Toggle Switch).

\subsection{Реализация системы запуска, мониторинга и отображения статусов CI/CD-пайплайнов}

\subsubsection{Асинхронные адаптеры интеграции}
Для взаимодействия с внешними системами автоматизации разработан сервисный уровень с базовым интерфейсом \texttt{BaseCIAdapter}. За счет этого архитектура приложения не зависит от конкретной платформы.

Реализация адаптера для GitHub Actions использует клиент \texttt{httpx.AsyncClient} с пулом соединений. Фрагмент кода отправки команды на повторный запуск конвейера (Re-run) приведен в листинге~\ref{lst:github_rerun}.

\begin{figure}[H]
\centering
\begin{minipage}{0.95\textwidth}
\footnotesize
\begin{verbatim}
import httpx
from fastapi import HTTPException

class GitHubAdapter:
    def __init__(
        self,
        token: str,
        base_url: str = "https://api.github.com"
    ):
        self.base_url = base_url
        self.headers = {
            "Authorization": f"Bearer {token}",
            "Accept": "application/vnd.github+json",
            "X-GitHub-Api-Version": "2022-11-28"
        }

    async def rerun_workflow(
        self,
        repo_full_name: str,
        run_id: str
    ) -> bool:
        url = (
            f"{self.base_url}/repos/{repo_full_name}"
            f"/actions/runs/{run_id}/rerun"
        )
        async with httpx.AsyncClient(timeout=10.0) as client:
            resp = await client.post(url, headers=self.headers)
            if resp.status_code == 201:
                return True
            raise HTTPException(
                status_code=resp.status_code,
                detail=f"Ошибка внешнего API: {resp.text}"
            )
\end{verbatim}
\end{minipage}
\caption{Асинхронный метод перезапуска пайплайна через GitHub Actions API}
\label{lst:github_rerun}
\end{figure}

\subsubsection{Отображение состояния сборок на клиентской части}
Главный экран мониторинга (\texttt{DashboardPage}) агрегирует данные обо всех активных пайплайнах. Статусы сборок нормализуются к единому формату перечисления (\texttt{QUEUED}, \texttt{RUNNING}, \texttt{SUCCESS}, \texttt{FAILED}, \texttt{CANCELED}) вне зависимости от исходного провайдера.

Каждая карточка содержит визуальные индикаторы:
\begin{itemize}
    \item зеленый индикатор --- конвейер выполнен успешно;
    \item красный пульсирующий бейдж --- зафиксирована ошибка выполнения стадии;
    \item синий анимированный индикатор --- пайплайн находится в процессе выполнения;
    \item контекстная кнопка действия «Re-run» --- активна для упавших или отмененных сборок. При нажатии вызывается мутация данных через REST API, а статус немедленно переходит в состояние ожидания запуска.
\end{itemize}

\subsection{Реализация системы логирования, хранения и обработки данных запусков}

\subsubsection{Сбор и нормализация журналов выполнения}
При обнаружении сбоя в работе конвейера система выполняет запрос к внешнему API для получения вывода консоли упавшего шага. Ввиду того, что полный объем журналов сборки может достигать десятков мегабайт, в базе данных PostgreSQL сохраняется только информативный срез (последние 200 строк вывода ошибки, извлекаемые регулярными выражениями по ключевым маркерам \texttt{ERROR}, \texttt{FATAL}, \texttt{FAILED}).

На фронтенде для отображения журналов разработан компонент \texttt{LogViewer}. Он использует моноширинный шрифт, темную цветовую схему и поддерживает автоматическую прокрутку к месту возникновения критической ошибки.

\subsubsection{Агрегация и визуализация аналитических метрик}
Для оценки эффективности процессов непрерывной интеграции реализован контроллер аналитики \texttt{/api/v1/metrics/summary}. Сервис агрегирует исторические записи из таблицы \texttt{pipelines} с применением математических функций PostgreSQL:
\begin{equation}
    \text{Success Rate} = \frac{N_{\text{success}}}{N_{\text{total}}} \times 100\%,
\end{equation}
где $N_{\text{success}}$ --- количество успешно завершенных сборок за выбранный период, $N_{\text{total}}$ --- общее число запусков конвейеров.

На клиентской стороне вычисленные показатели отображаются в виде обзорных карточек метрик (KPI widgets) и временных графиков, позволяющих отслеживать среднюю продолжительность сборки и динамику процента падений по дням недели.

Разработанный программный комплекс полностью реализует поставленные функциональные требования, формируя надежный и удобный инструмент централизованного мониторинга сборок.